\section{Discussion}
\label{sec:discussion}

\subsection{Ethics and Accountability}
\label{sec:ethics}

The system leaves the final decision to a human reviewer. That is a deliberate design
choice. Automatic rejection would save time, but it would also hide responsibility inside the
pipeline; in human-in-the-loop review, responsibility must land on a specific person and a
specific decision. The audit trail is therefore not an add-on but part of the system: it
records the model proposal, the reviewer's action, the time and the reason, so that the
decision can be reviewed afterwards.

We follow data minimisation: the evaluation uses synthetic documents only. A real deployment
should first settle data authorisation, de-identification and retention limits. Model cards
and datasheets should be maintained with the system \cite{mitchell2019,gebru2021}, not
written only when an audit arrives.

\subsection{Limitations}
\label{sec:limitations}

The evaluation has three clear limits. First, the corpus is synthetic and small, with only
two topics, so the time saving may not generalise to real documents. Second, there are only
two reviewers, and reviewer differences are not analysed separately. Third, the evaluation
does not cover automation bias over time: reviewers may gradually reduce their independent
judgement. The next step needs a larger and more varied corpus, and a longitudinal
observation period.
